% main.tex -- Paper 2: what is the GPU speedup measured against?
%
% Secondary analysis of the pre-registered four-host study reported in Paper 1
% (../paper/). Every number is a \stat{} / \ratioci{} / \ratiotext{} macro filled
% from the analysis outputs by tools/stats_to_tex.py; nothing is typed in.
% Statements made in words about the data ("at every scale", "every interval
% below 1") are checked by the converter's claim list and marked here with the
% invisible \claim{<id>}; a claim that stops holding fails the build.
%
% Conditional text (tools/stats_to_tex.py, macros.tex): regions between
% "% BEGIN-IF:<cond>" and "% END-IF:<cond>" (IFNOT for the else branch) inside
% the matching \if<flag> conditional. Conditions: bridge_a100, bridge_h100,
% bridge_gh200, bridge_b200 (the host's side study and E1/E4 are present),
% bridge_extra (A100 or B200 bridge present), pcheck (protocol_check.* keys),
% decomp (decomp.* keys).
%
% Labelling rule: every post-hoc number carries the word "exploratory" in the
% sentence, caption or table header that reports it, and \expl in tables.
\documentclass[11pt]{article}

\usepackage[T1]{fontenc}
\usepackage[utf8]{inputenc}
\usepackage{lmodern}
\usepackage[margin=1in]{geometry}
\usepackage{amsmath,amssymb}
\usepackage{booktabs}
\usepackage{graphicx}
\usepackage[table]{xcolor}
\usepackage[numbers,sort&compress]{natbib}
\usepackage{caption}
\usepackage{subcaption}
\usepackage{enumitem}
\usepackage{xspace}
\usepackage{array}
\usepackage[hidelinks]{hyperref}
\usepackage{url}
\usepackage{authblk}
\usepackage{seqsplit}

% macros.tex (Paper 2) -- statistics lookup, figure inclusion and notice macros.
%
% Numbers in the paper are never typed in. Every result is written as \stat{key},
% where key is a dotted stats.json key (e.g. c1.baseline.sf300.estimate).
% paper2/tools/stats_to_tex.py turns the host stats.json files into
% generated/stats.tex, which contains one \defstat{key}{value} per key and
% \statsfinaltrue. In final mode a missing key is a hard LaTeX error; in draft
% mode (no generated/stats.tex) the key is printed in a box so the gaps are visible.

\makeatletter
\newif\ifstatsfinal
\statsfinalfalse
\newcommand{\defstat}[2]{\expandafter\def\csname stat@#1\endcsname{#2}}
\newcommand{\statplaceholder}[1]{%
  \fcolorbox{red}{yellow!20}{\scriptsize\ttfamily\detokenize{#1}}}
\newcommand{\stat}[1]{%
  \ifcsname stat@#1\endcsname
    \csname stat@#1\endcsname
  \else
    \ifstatsfinal
      \PackageError{stats}{Missing statistic '#1' in generated/stats.tex}%
        {Re-run tools/stats_to_tex.py; the analysis output lacks this key.}%
    \else
      \statplaceholder{#1}%
    \fi
  \fi}
% "estimate [ci_low, ci_high]" for a ratio prefix such as c1.baseline.sf300
\newcommand{\ratioci}[1]{%
  \stat{#1.estimate}\,{\footnotesize[\stat{#1.ci_low},\,\stat{#1.ci_high}]}}
% "estimate x (95% CI ci_low--ci_high)" for running text
\newcommand{\ratiotext}[1]{%
  \stat{#1.estimate}$\times$ (95\%~CI \stat{#1.ci_low}--\stat{#1.ci_high})}
\makeatother

% Conditions (paper2/tools/stats_to_tex.py). Each is true in a draft so every
% branch shows its placeholders; the converter writes \<flag>false when the input
% is absent. Wrap dependent text in \if<flag> ... \fi between
% "% BEGIN-IF:<cond>" / "% END-IF:<cond>" comment lines (IFNOT for the \else
% branch) so the key scanner knows which keys each case needs.
\newif\ifbridgeaonehundred \bridgeaonehundredtrue  % bridge_a100: A100 side study + E1/E4
\newif\ifbridgehonehundred \bridgehonehundredtrue  % bridge_h100
\newif\ifbridgeghtwohundred \bridgeghtwohundredtrue % bridge_gh200
\newif\ifbridgebtwohundred \bridgebtwohundredtrue  % bridge_b200
\newif\ifbridgeextra \bridgeextratrue              % bridge_extra: A100 or B200 bridge
\newif\ifpcheck \pchecktrue                        % pcheck: protocol_check.* keys
\newif\ifdecomp \decomptrue                        % decomp: decomp.* keys
\newif\ifeadj \eadjtrue                            % e1adj: exploratory_adj.e1.* keys

% Disclosure watermarks, set by generated/stats.tex: hand-filled _display values
% (stats_to_tex.py --allow-manual-overrides, or a draft) and synthetic stats.
% Printed diagonally on every page and in the title block.
\newif\ifmanualvalues
\newif\ifsyntheticstats

\IfFileExists{generated/stats.tex}{\input{generated/stats.tex}}{}

\newcommand{\watermarktext}{%
  \ifmanualvalues MANUAL VALUES\fi
  \ifmanualvalues\ifsyntheticstats\ \textperiodcentered\ \fi\fi
  \ifsyntheticstats SYNTHETIC -- NOT RESULTS\fi}
\newcommand{\watermarkpage}{%
  \put(0.5\paperwidth,-0.5\paperheight){\makebox(0,0){\rotatebox{45}{%
    \scalebox{4}{\sffamily\bfseries\textcolor{red!30}{\watermarktext}}}}}}
\AddToHook{shipout/background}{%
  \ifmanualvalues\watermarkpage\else\ifsyntheticstats\watermarkpage\fi\fi}
\newcommand{\watermarknotice}{%
  \ifmanualvalues\\\textcolor{red}{\textbf{\watermarktext}}%
  \else\ifsyntheticstats\\\textcolor{red}{\textbf{\watermarktext}}\fi\fi}

% Figures are copied into figures/ by tools/sync_figures.py (never drawn here). Until then a
% labelled placeholder box keeps the layout and cross-references working.
\newcommand{\studyfig}[2]{%
  \IfFileExists{figures/#1.pdf}{\includegraphics[width=#2]{figures/#1.pdf}}{%
  \IfFileExists{figures/#1.png}{\includegraphics[width=#2]{figures/#1.png}}{%
  \fbox{\parbox[c][0.28\textheight][c]{0.9\linewidth}{\centering\ttfamily
    figure pending: \detokenize{#1}}}}}}
% \studyfigalt{stem}{fallback-stem}{width}: the first figure that exists.
\newcommand{\studyfigalt}[3]{%
  \IfFileExists{figures/#1.pdf}{\studyfig{#1}{#3}}{%
  \IfFileExists{figures/#1.png}{\studyfig{#1}{#3}}{\studyfig{#2}{#3}}}}
\newcommand{\iffigure}[3]{\IfFileExists{figures/#1.pdf}{#2}{\IfFileExists{figures/#1.png}{#2}{#3}}}

% Author-supplied detail still missing (affiliation, email, venue, ...).
\newcommand{\todoplaceholder}[1]{\textcolor{red}{[TODO: #1]}}

% Results-not-yet-available marker for draft prose (also flagged with
% "% TODO(results)" comments in the source).
\newcommand{\pending}[1]{\textcolor{gray}{[pending results: #1]}}

% TPC fair-use disclaimer. The template wording is TPC's (Fair Use of TPC
% Benchmarks quick reference, 2019-06-05; current Policies v6.20 sec. 8.1.5).
% The name of the derived work is an open author decision; see paper/README.md.
\newcommand{\workloadname}{H100 study workload}
\newcommand{\tpcdisclaimer}{The \workloadname{} is derived from the TPC-H
  Benchmark and as such is not comparable to published TPC-H Benchmark results,
  as the \workloadname{} results do not comply with the TPC-H Benchmark
  Specification.}
\newcommand{\chartnote}{\emph{Derived from TPC-H; not comparable to published
  TPC-H results (see Section~\ref{sec:scope}).}}

% Engine and contrast names
\newcommand{\eng}[1]{\texttt{#1}}
\newcommand{\polarscpu}{\eng{polars}}
\newcommand{\polarsgpu}{\eng{polars\_gpu}}
\newcommand{\sparkcpu}{\eng{spark\_cpu}}
\newcommand{\sparkgpu}{\eng{spark\_gpu}}
\newcommand{\q}[1]{\texttt{#1}}

% \claim{id}: invisible marker for a statement the converter re-checks against the data
% (tools/stats_to_tex.py, claims()); an id with no check is a build error.
\newcommand{\claim}[1]{}

% Paper 2 shorthands
\newcommand{\polarsold}{\stat{host.packages.polars}}               % Polars under cudf-polars
\newcommand{\polarsnew}{\stat{exploratory.e1.meta.cpu_version}}    % Polars in the bridge
\newcommand{\expl}{\textsuperscript{\textsc{x}}}                    % marks an exploratory value

\setlength{\emergencystretch}{3em}

% TODO(author): title options.
%  (a) [current] What Is the GPU Speedup Measured Against? Library-Version
%      Lock-In and Core-Count Confounds in CPU-versus-GPU Dataframe Benchmarks
%  (b) Same Threads, Different Baselines: How the CPU Side Shapes GPU Speedups
%      for Polars and Spark
%  (c) The Denominator Problem in GPU Dataframe Benchmarks
% The subtitle states the secondary and partly exploratory nature; keep it in
% any variant.
\title{What Is the GPU Speedup Measured Against?\\
  Library-Version Lock-In and Core-Count Confounds\\
  in CPU-versus-GPU Dataframe Benchmarks\\[4pt]
  \large A secondary analysis of a pre-registered four-host study;\\
  the re-baselined and cross-host analyses are exploratory}
\author{Daniel Ejiofor}
\affil{Boston University, Boston, MA, USA}
\affil{\href{mailto:ejiod@bu.edu}{\texttt{ejiod@bu.edu}}}
\date{Draft of \today\watermarknotice}

\begin{document}
\maketitle

\begin{abstract}
A GPU speedup is a ratio, so what it means depends on its CPU denominator. We re-examine the CPU baseline of a pre-registered benchmark of Polars and Apache
Spark, each on CPU and GPU, run as four separate campaigns on single-GPU A100,
H100, GH200 and B200 hosts with a TPC-H-derived workload (not a TPC-H
result)~\citep{ejiofor2026companion}.

\emph{Version lock-in.} The Polars GPU engine requires Polars below 1.45, so
its pre-registered speedup is measured against Polars \polarsold{} on the CPU.
In the pre-registered bridge, Polars \polarsnew{} took
\stat{p2.bridge.baseline.min}--\stat{p2.bridge.baseline.max} of
\polarsold's time on the baseline suite and
\stat{p2.bridge.complex.min}--\stat{p2.bridge.complex.max} on the complex
suite. Every 95\% interval lay below~1
\stat{p2.bridge.scope}\stat{p2.bridge.exception_short}\claim{bridge-never-slower}.
Measured against \polarsnew{} instead (exploratory), the H100 speedup at SF300
falls from \ratiotext{c1.baseline.sf300} to
\ratiotext{exploratory.e1.baseline.sf300}, or to
\ratiotext{exploratory_adj.e1.h100.baseline.sf300} after correcting for a
measured session bias. \stat{p2.reversal.phrase_cap}, the GPU engine is faster
than Polars \polarsold{} but slower than Polars
\polarsnew{}\claim{reversal-exists}.

\emph{Equal threads, unequal cores.} The same \stat{host.cpus_arg} threads occupy
\stat{p2.h100.cores_used} physical cores on the H100 and B200 (two hardware
threads per core) but \stat{p2.a100.cores_used} on the A100 and
GH200\claim{cores-pairs}. The GH200's Polars GPU speedup at SF300 was
\stat{gh200.c1.baseline.sf300.estimate}$\times$, against
\stat{c1.baseline.sf300.estimate}$\times$ on the H100, mainly because its CPU
side is faster (exploratory
factorization)\claim{gh200-cpu-side}. The CPU model and architecture also
differ, so core count is not isolated.

We end with a reporting checklist for CPU-versus-GPU benchmarks.
\end{abstract}

\noindent\textbf{Scope notice.} \tpcdisclaimer{} The queries are a selected
and adapted subset with additional custom queries, Float64 normalization of
decimal columns, and our own timing and summary method; see
Section~\ref{sec:scope}.

% ======================================================================
\section{Introduction}
\label{sec:intro}

Claims that a GPU accelerates analytical queries are usually reported as one
number: the time on a CPU engine divided by the time on the GPU engine. Such a
ratio has two sides. Most of the attention, in benchmarks and in the
literature, goes to the GPU side: the device, its memory, the operators it
supports. The CPU side is often treated as a fixed reference. Earlier work has
shown that this reference matters. Lee et~al.\ found that reported
GPU advantages of 10 to 1000 times on throughput kernels narrowed to 2.5 times
on average once both the CPU and the GPU code were
optimized~\citep{lee2010debunking}. Shanbhag et~al.\ argued that, for
memory-bound analytical operators, the gain to expect is about the ratio of GPU
to CPU memory bandwidth rather than the orders of magnitude previously claimed,
and measured where full queries depart from it~\citep{shanbhag2020gpucpu}. Database benchmarking
guides list unequal tuning and unreported configuration among the commonest
pitfalls~\citep{raasveldt2018fair,hoefler2015scientific}.

This paper examines two such issues in a setting where they arise from the
software and hardware themselves, not from the experimenter's effort. The data
come from a completed, pre-registered study~\citep{ejiofor2026companion}
(Paper~1). It compares Polars on the CPU, Polars with its cuDF-based GPU
engine, Spark on the CPU and Spark with NVIDIA's cuDF plugin, run as four
separate campaigns with an identical design on single-GPU A100, H100, GH200
and B200 hosts. Paper~1 reports the pre-registered contrasts. Here we ask what
those contrasts are measured \emph{against}.

\begin{enumerate}[leftmargin=*,itemsep=2pt]
  \item \textbf{Version lock-in.} The Polars GPU engine ships as a separate
    package that supports only a range of Polars versions. When the study was
    locked, the newest release required Polars below
    1.45~\citep{pypi_cudfpolars_26102}, while Polars~2.0.0 and~2.1.0 had already
    been released~\citep{pypi_polars}. The pre-registered within-framework
    speedup therefore compares the GPU engine with an older CPU release. A
    side study (pre-registered as RQ5) measured how much the CPU release
    matters. We use it to ask how the GPU speedup changes when it is expressed
    against the newer CPU release. That second step is post hoc.
  \item \textbf{Equal threads are not equal cores.} To keep the hosts
    comparable, every engine on every host ran with the same number of
    worker threads. The hosts' virtual machines expose different CPU
    topologies, however. On two hosts each core carries two hardware threads;
    on the other two it carries one. The same thread count therefore sits on
    half as many physical cores on the first pair. We ask how much of the
    cross-host difference in GPU speedup comes from the CPU side of the ratio.
\end{enumerate}

\paragraph{Contributions.}
(i)~A re-baselining of the Polars GPU speedup against the newer Polars CPU
release on every host with a version bridge, including where its verdict
changes (Section~\ref{sec:res-version}). This analysis is exploratory.
(ii)~A factorization of cross-host differences in GPU speedup into CPU-side and
GPU-side time ratios, which holds exactly for the point estimates when query
sets match. Applied to the four hosts, it shows that the topology recorded on
each host lines up with the CPU-side factor
(Section~\ref{sec:res-cores}). This analysis is exploratory and descriptive.
(iii)~A case study of how an earlier, non-pre-registered A100 study's baseline
choices carry into a comparison with this one (Section~\ref{sec:res-prior}).
(iv)~A reporting checklist for CPU-versus-GPU dataframe benchmarks
(Section~\ref{sec:checklist}).

\paragraph{What this paper is not.} It is a secondary analysis. It reuses the
primary study's data, adds no new measurements, and does not restate the
primary results; for those, see Paper~1. Every analysis below is labelled
either as part of the locked design or as exploratory
(Table~\ref{tab:status}). Exploratory results generate hypotheses; they are
not confirmatory evidence~\citep{nosek2018prereg,gelman2014crisis}.

% ======================================================================
\section{Background}
\label{sec:background}

\paragraph{Engines.} Polars is a columnar dataframe library with a streaming
CPU engine~\citep{polars1442,polars200}. Since 2024 it can hand a query plan
to a GPU engine built on NVIDIA cuDF, distributed as the separate package
\texttt{cudf-polars}~\citep{semaan2024polarsgpu,polars_gpu_guide,cudf_polars_docs}.
Spark SQL~\citep{armbrust2015sparksql} runs on the GPU through NVIDIA's cuDF
plugin (formerly the RAPIDS Accelerator)~\citep{nvidia_cudf_spark_docs}. GPU
query processing has a long history, surveyed
in~\citep{bress2014gpudb,rosenfeld2022heterogeneous}. On machines like these,
the interconnect and memory system often decide the
outcome~\citep{lutz2020pump,shanbhag2020gpucpu}.

\paragraph{The release coupling.} \texttt{cudf-polars} works on Polars' internal
plan, so each release declares the range of Polars versions it supports. The
release used here, 26.10.2, requires Polars ${\ge}1.35,{<}1.45$ and Python
${\ge}3.11$~\citep{pypi_cudfpolars_26102}. Polars~1.44.2 is the last 1.x
release; 2.0.0 and 2.1.0 followed on 2026-10-06 and
2026-10-25~\citep{pypi_polars}. A GPU-engine user therefore runs a Polars
release older than the newest one. A benchmark that uses one Polars install
for both engines inherits that older release for its CPU baseline.

\paragraph{Simultaneous multithreading.} With simultaneous multithreading (SMT),
one physical core runs instructions from more than one hardware thread, which
share the core's execution units and caches~\citep{tullsen1995smt}. A thread
count therefore does not say how much core hardware it uses: 26 threads on 13
two-way SMT cores share half the execution resources of 26 threads on 26
cores. In cloud virtual machines, the guest operating system reports the
topology the hypervisor exposes. That topology need not match the physical
machine.

\paragraph{Fair comparison and pre-registration.} Benchmarking guides ask
authors to report and justify every configuration choice that could change a
result~\citep{raasveldt2018fair,hoefler2015scientific,kalibera2013rigorous}.
Even setup details that look harmless can bias
measurements~\citep{mytkowicz2009wrong}. Pre-registration separates analyses
that were fixed before the data existed from those chosen after seeing
them~\citep{nosek2018prereg}. Results reached through many possible analysis
choices are more likely to be chance findings~\citep{gelman2014crisis}.

% ======================================================================
\section{Data and Methods}
\label{sec:methods}

\subsection{Source study}
\label{sec:source}

Paper~1~\citep{ejiofor2026companion} describes the design in full; we
summarize what this analysis depends on. The design was locked before any
timed data existed. It was committed as \texttt{ee2e6a7}, and amendment A4 as
\texttt{afd07c0}, in the study repository. The H100 campaign ran under design
version \texttt{\stat{meta.design_commit_short}}, the GH200 campaign under
\texttt{\stat{gh200.meta.design_commit_short}}, the A100 campaign under
\texttt{\stat{a100.meta.design_commit_short}} and the B200 campaign under
\texttt{\stat{b200.meta.design_commit_short}}. Amendments A2 and A3 give each
host its own complete campaign with the same seed, order design, gates and side
study, and data are never pooled across hosts. The H100 is the primary host.

\begin{description}[leftmargin=*,itemsep=2pt]
  \item[Arms.] \polarscpu{} (Polars \polarsold, streaming engine),
    \polarsgpu{} (\texttt{cudf-polars} \stat{host.packages.cudf-polars-cu12} on
    the same Polars install), \sparkcpu{} (Spark~3.5.9, \texttt{local[26]}) and
    \sparkgpu{} (Spark with the cuDF plugin). Settings are fixed and untuned.
    On every host, every engine ran with \texttt{-{}-cpus}~\stat{host.cpus_arg}.
    This sets the engines' worker pools (\texttt{POLARS\_MAX\_THREADS},
    \texttt{local[N]}). It does not pin threads to cores or cap the total number
    of operating-system threads, and the code sets no CPU affinity.
  \item[Workload.] A baseline suite of six adapted TPC-H queries at SF1--SF300,
    and a complex suite of six queries (two TPC-H-derived, four custom) at
    SF30--SF300, on Parquet inputs regenerated and verified on each host.
    Polars GPU cannot run two complex queries, \q{cx\_q13} and
    \q{cx\_window}; this gap was pre-declared. Ratios involving Polars GPU
    exclude these queries.
  \item[Protocol.] Four rounds per host, ordered by a Williams Latin
    square~\citep{williams1949latin}. One session per engine and (suite, scale)
    per round: preheat, one warm-up, then three measured trials per query. The
    unit is the median of the three trials. Output-agreement, GPU-residency and
    cleanup gates never depend on timings.
  \item[Contrasts.] Speedups $t_{\mathrm{ref}}/t_{\mathrm{engine}}$ per query
    and round. Suite values are geometric means~\citep{fleming1986geomean} over
    the queries supported on both sides, computed per round, with a 95\%
    $t$-interval across rounds (df~$=3$). This paper uses C1
    (\polarscpu{}/\polarsgpu{}), C2 (\sparkcpu{}/\sparkgpu{}) and C4
    (\polarscpu{}/the faster GPU engine per query and round).
\end{description}

\subsection{What was pre-registered and what was not}
\label{sec:status}

Table~\ref{tab:status} classifies every analysis in this paper. The
classification follows the locked design document, not our view of how
convincing a result is.

\begin{table}[t]
\centering
\caption{Pre-registration status of each analysis in this paper.}
\label{tab:status}
\small
\begin{tabular}{@{}p{5.2cm}p{3.1cm}p{6.6cm}@{}}
\toprule
Analysis & Status & Source and notes \\
\midrule
C1, C2, C4 per host & pre-registered & design \S9; reported in Paper~1 \\
Version bridge $t_{\text{\polarsnew}}/t_{\text{\polarsold}}$ (RQ5) & pre-registered & design \S10, on every host by A2/A3; paired $t$-interval, df~$=3$ \\
Cross-host time ratios $t_{\text{host}}/t_{\text{H100}}$ & pre-registered as descriptive & A2 secondary analysis; Welch-type intervals; CPU-engine ratios ``only with the caveat'' that CPU and RAM differ \\
Host identity (GPU, CPU model, RAM, driver) & pre-registered recording & A2: recorded in each campaign \\
CPU topology fields (SMT, cores, NUMA) & recorded as part of the pre-registered inventory; analysed post hoc & the runner's standard inventory saves the guest's \texttt{lscpu} output in every session; the design (\S6, \S9) does not name topology as a factor \\
E1, E4 (re-baselined speedups) & \textbf{exploratory} & added after the bridge result was seen; Welch-type intervals \\
First-above-1 scale for E1/E4 & \textbf{exploratory} & same rule as the pre-registered RQ3 crossover, applied to E1/E4 \\
Speedup factorization across hosts & \textbf{exploratory} & Section~\ref{sec:factor}; descriptive \\
Bridge over GPU-supported queries & \textbf{exploratory} & point estimates only \\
% BEGIN-IF:pcheck
\ifpcheck
Protocol check K1 (side-study vs main-round Polars~\polarsold) & \textbf{exploratory} & Section~\ref{sec:pcheck}; Welch-type intervals \\
Session-adjusted E1 (E1/K1) & \textbf{exploratory} & Section~\ref{sec:pcheck}; point estimates, and intervals where the analysis provides them \\
\fi
% END-IF:pcheck
Comparison with the earlier A100 study & descriptive, illustrative & the earlier study was not pre-registered (no locked design; one session per cell) \\
\bottomrule
\end{tabular}
\end{table}

\subsection{Version bridge (pre-registered) and re-baselined speedups (exploratory)}
\label{sec:rebase}

After its main rounds, each host's side study compares Polars \polarsold{}
(main environment) with Polars \polarsnew{} (a separate environment) on the
CPU. It covers both suites and all scales, in four rounds ordered AB, BA, BA,
AB, with the same session protocol. The bridge estimate is the suite
geometric-mean ratio $B = t_{\text{\polarsnew}}/t_{\text{\polarsold}}$ with a
paired $t$-interval (df~$=3$). Below~1 means \polarsnew{} is faster.

The exploratory ratios E1 and E4 replace the numerator of C1 and C4 with the
side study's Polars \polarsnew{} times:
\[
  \mathrm{E1} = \frac{t_{\text{Polars \polarsnew{} CPU}}}{t_{\text{Polars GPU}}},
  \qquad
  \mathrm{E4} = \frac{t_{\text{Polars \polarsnew{} CPU}}}{t_{\text{best GPU engine}}}.
\]
They use the same queries as their pre-registered counterparts. The CPU times
come from side-study sessions and the GPU times from main-round sessions, so
the two sides are unpaired; the intervals are Welch-type on the log
scale~\citep{welch1947,satterthwaite1946}. In this build each is a
\stat{exploratory_summary.interval_note}. Two consequences follow. E1 and C1
differ in the CPU release, but also in which sessions supplied the CPU times.
And an E1 interval is not comparable in width to a C1 interval.
% BEGIN-IF:pcheck
\ifpcheck
Section~\ref{sec:pcheck} measures the session effect directly. A shift from C1
to E1 counts as distinguishable from it only where the shift exceeds the K1
interval's largest distance from~1. It defines
$\mathrm{K1} = t^{\text{side}}_{\text{\polarsold}}/t^{\text{main}}_{\text{\polarsold}}$,
the same release timed in the two kinds of session, and the session-adjusted
ratio $\mathrm{E1}/\mathrm{K1}$. For this ratio K1 is taken over E1's queries. On the complex suite these are
the four GPU-supported queries, whereas the K1 values reported in
Section~\ref{sec:pcheck} use all six Polars queries. With that choice,
$\mathrm{E1}/\mathrm{K1}$ equals the bridge over E1's queries times C1, and both
factors are paired within their own rounds.
\fi
% END-IF:pcheck
Where we say that the GPU side ``wins'' or ``loses'', the whole 95\% interval
lies above or below~1. We apply the pre-registered RQ3 rule (the first scale at
which the interval lies entirely above~1) to E1 and E4 only as a description.
In every case reported here, all larger scales also lie above~1\claim{first-above-sustained}.

\subsection{CPU topology}
\label{sec:topology-method}

The runner's standard inventory, which predates this study, saved the guest's
\texttt{lscpu} output in every session's environment record. Recording it was
part of the pre-registered inventory; analysing CPU topology was not
pre-registered, and the design does not name topology as a factor (\S6,
\S9). From it the analysis
takes the hardware threads per core, cores per socket, sockets, NUMA nodes
(with and without CPUs) and the hypervisor line. It derives the range of
physical cores that \texttt{-{}-cpus} threads can occupy: at least
$\lceil \min(\text{threads},\text{vCPUs})/\text{SMT}\rceil$ when threads pack
onto sibling hardware threads, and at most
$\min(\text{threads},\text{vCPUs},\text{cores})$ when they spread out. These
values are \emph{guest-reported}. A hypervisor can expose a topology that
differs from the physical machine, and lacking a hypervisor line does not
prove bare metal.

\subsection{Factorizing a cross-host speedup difference}
\label{sec:factor}

Let $C^{h}$ be a host's suite speedup $t_{\text{cpu}}/t_{\text{gpu}}$ and
$X^{h}_{e}$ the cross-host time ratio $t^{h}_{e}/t^{\text{H100}}_{e}$ of engine
$e$. Both are geometric means over the same queries and rounds. Taking logs and
reordering the means gives, for the point estimates,
\begin{equation}
  \frac{C^{h}}{C^{\text{H100}}} \;=\; \frac{X^{h}_{\text{cpu}}}{X^{h}_{\text{gpu}}}.
  \label{eq:factor}
\end{equation}
A host's speedup can thus be lower than the H100's for two reasons: its CPU
engine is faster ($X_{\text{cpu}}<1$) or its GPU engine is slower
($X_{\text{gpu}}>1$). The identity is exact when the speedup and both
cross-host ratios use the same queries. In our data that holds for every cell
except C1 on the complex suite\claim{speedup-identity}\claim{identity-only-c1-complex}. There, the cross-host Polars CPU ratio
includes the two queries the GPU engine cannot run.
% BEGIN-IF:decomp
\ifdecomp
For C1 on the complex suite we therefore use a CPU-side ratio restricted to
the GPU-supported queries, and an interval for $C^{h}/C^{\text{H100}}$ from the
per-round log speedups of the two hosts (Welch-type)\claim{decomp-identity}.
\fi
% END-IF:decomp
% BEGIN-IFNOT:decomp
\ifdecomp\else
We report the factorization only where the identity is exact.
\fi
% END-IFNOT:decomp
Equation~\eqref{eq:factor} is arithmetic, not a causal decomposition. A CPU-side
factor below~1 says the host's CPU engine was faster. It does not say why: core
count, microarchitecture, memory system and architecture all change together
across these hosts.

\subsection{Scope and notice}
\label{sec:scope}
\tpcdisclaimer{} Deviations from the TPC-H specification include: a selected
subset of six adapted queries with fixed substitution parameters, plus two
TPC-H-derived and four custom complex queries; data from DuckDB's generator,
which differs slightly from the specification~\citep{duckdb_tpch_ext};
Float64 instead of decimal arithmetic; extra sort keys; Parquet files rather
than a loaded database; no refresh functions, power or throughput tests, or
TPC metrics; and our own timing boundary and summary statistics. No result in
this paper may be compared with official TPC-H
results~\citep{tpch301,boncz2014tpch}. Use of the benchmark name and the TPC
software is subject to the TPC
policies~\citep{tpcpolicies620,tpcfairuse2019} and the TPC
EULA~\citep{tpceula22}.
% TODO(author): same open TPC permission question as Paper 1 (paper/README.md).

% ======================================================================
\section{Results}
\label{sec:results}

\subsection{Version lock-in}
\label{sec:res-version}

\paragraph{The bridge (pre-registered).}
Table~\ref{tab:bridge} and Figure~\ref{fig:bridge} give the version bridge \stat{p2.bridge.scope}. The bridge interval lies below~1 \stat{p2.bridge.below_phrase}. No
interval lies above~1, so Polars \polarsnew{} was never slower than
\polarsold{}\claim{bridge-never-slower}. On the
baseline suite the ratio ranged from \stat{p2.bridge.baseline.min} to
\stat{p2.bridge.baseline.max}. On the H100 the estimate \stat{bridge.baseline.summary.trend}. On the complex suite it ranged from
\stat{p2.bridge.complex.min} to \stat{p2.bridge.complex.max}. One complex query
drives much of that: \q{cx\_window} ran in \stat{p2.bridge.cx_window.min} to
\stat{p2.bridge.cx_window.max} of its \polarsold{} time under \polarsnew{}.
Polars GPU cannot run that query, so it does not enter the GPU contrasts.
Restricted to the queries that do enter them (exploratory; point estimates),
the complex-suite bridge is between \stat{p2.bridge_gpuq.complex.min} and
\stat{p2.bridge_gpuq.complex.max}. A bridge quoted over the whole suite would
therefore overstate the version effect that matters for the GPU comparison.

\begin{table}[t]
\centering
\caption{Version bridge (pre-registered, RQ5): suite geometric-mean ratio
$t_{\text{\polarsnew}}/t_{\text{\polarsold}}$ on the CPU, with 95\%
$t$-intervals across four rounds. Below~1: \polarsnew{} is faster. The last
column (\expl{}, exploratory; point estimate) restricts the SF300 value to the
queries Polars GPU runs. \chartnote}
\label{tab:bridge}
\scriptsize
\setlength{\tabcolsep}{3pt}
\begin{tabular}{@{}llcccccc@{}}
\toprule
Host & Suite & SF1 & SF10 & SF30 & SF100 & SF300 & \shortstack{SF300,\\GPU queries\expl} \\
\midrule
% BEGIN-IF:bridge_a100
\ifbridgeaonehundred
A100 & baseline & \ratioci{a100.bridge.baseline.sf1} & \ratioci{a100.bridge.baseline.sf10} & \ratioci{a100.bridge.baseline.sf30} & \ratioci{a100.bridge.baseline.sf100} & \ratioci{a100.bridge.baseline.sf300} & \stat{p2.a100.bridge_gpuq.baseline.sf300} \\
     & complex  & --- & --- & \ratioci{a100.bridge.complex.sf30} & \ratioci{a100.bridge.complex.sf100} & \ratioci{a100.bridge.complex.sf300} & \stat{p2.a100.bridge_gpuq.complex.sf300} \\
\midrule
\fi
% END-IF:bridge_a100
H100 & baseline & \ratioci{bridge.baseline.sf1} & \ratioci{bridge.baseline.sf10} & \ratioci{bridge.baseline.sf30} & \ratioci{bridge.baseline.sf100} & \ratioci{bridge.baseline.sf300} & \stat{p2.h100.bridge_gpuq.baseline.sf300} \\
     & complex  & --- & --- & \ratioci{bridge.complex.sf30} & \ratioci{bridge.complex.sf100} & \ratioci{bridge.complex.sf300} & \stat{p2.h100.bridge_gpuq.complex.sf300} \\
\midrule
GH200 & baseline & \ratioci{gh200.bridge.baseline.sf1} & \ratioci{gh200.bridge.baseline.sf10} & \ratioci{gh200.bridge.baseline.sf30} & \ratioci{gh200.bridge.baseline.sf100} & \ratioci{gh200.bridge.baseline.sf300} & \stat{p2.gh200.bridge_gpuq.baseline.sf300} \\
      & complex  & --- & --- & \ratioci{gh200.bridge.complex.sf30} & \ratioci{gh200.bridge.complex.sf100} & \ratioci{gh200.bridge.complex.sf300} & \stat{p2.gh200.bridge_gpuq.complex.sf300} \\
% BEGIN-IF:bridge_b200
\ifbridgebtwohundred
\midrule
B200 & baseline & \ratioci{b200.bridge.baseline.sf1} & \ratioci{b200.bridge.baseline.sf10} & \ratioci{b200.bridge.baseline.sf30} & \ratioci{b200.bridge.baseline.sf100} & \ratioci{b200.bridge.baseline.sf300} & \stat{p2.b200.bridge_gpuq.baseline.sf300} \\
     & complex  & --- & --- & \ratioci{b200.bridge.complex.sf30} & \ratioci{b200.bridge.complex.sf100} & \ratioci{b200.bridge.complex.sf300} & \stat{p2.b200.bridge_gpuq.complex.sf300} \\
\fi
% END-IF:bridge_b200
\bottomrule
\end{tabular}
\end{table}

\begin{figure}[t]
\centering
\begin{subfigure}[t]{0.49\linewidth}\centering
\studyfig{f08_bridge_h100}{\linewidth}\caption{H100}\end{subfigure}\hfill
\begin{subfigure}[t]{0.49\linewidth}\centering
\studyfig{f08_bridge_gh200}{\linewidth}\caption{GH200}\end{subfigure}
% BEGIN-IF:bridge_extra
\ifbridgeextra
\par\medskip
% BEGIN-IF:bridge_a100
\ifbridgeaonehundred
\begin{subfigure}[t]{0.49\linewidth}\centering
\studyfig{f08_bridge_a100}{\linewidth}\caption{A100}\end{subfigure}\hfill
\fi
% END-IF:bridge_a100
% BEGIN-IF:bridge_b200
\ifbridgebtwohundred
\begin{subfigure}[t]{0.49\linewidth}\centering
\studyfig{f08_bridge_b200}{\linewidth}\caption{B200}\end{subfigure}
\fi
% END-IF:bridge_b200
\fi
% END-IF:bridge_extra
\caption{Version bridge (pre-registered, RQ5) per host:
$t_{\text{\polarsnew}}/t_{\text{\polarsold}}$ per suite and scale with 95\%
intervals. Below~1: \polarsnew{} is faster. Figures from the study's analysis
(F8). \chartnote}
\label{fig:bridge}
\end{figure}

% BEGIN-IFNOT:bridge_a100
\ifbridgeaonehundred\else
The A100 side study was not part of the analysis inputs for this build, so the
A100 appears in the cross-host analyses (Section~\ref{sec:res-cores}) but not
here.
\fi
% END-IFNOT:bridge_a100
% BEGIN-IFNOT:bridge_b200
\ifbridgebtwohundred\else
The same holds for the B200 side study.
\fi
% END-IFNOT:bridge_b200

\paragraph{Re-baselined speedups (exploratory).}
Table~\ref{tab:rebase} sets E1 and E4 beside C1 and C4, and
Figure~\ref{fig:sensitivity} plots them against scale. \stat{p2.bridge.scope_cap}, E1 is below C1 and E4 below C4 (point estimates)
\stat{p2.e1.below_phrase_ref}\claim{e1-exceptions-explained}\claim{e4-exceptions-explained}\claim{e1-e4-same-exceptions}. This is what the
bridge implies, because the GPU side is identical. On the H100, measuring
against \polarsnew{} instead of \polarsold{} lowers the Polars GPU speedup by
\stat{p2.h100.gap.baseline.range} on the
baseline suite, depending on scale, and by
\stat{p2.h100.gap.complex.range} on the complex
suite. The reductions are: \stat{p2.gap.by_host}
(all $100(1-\mathrm{E1}/\mathrm{C1})$, point estimates). In
\stat{p2.k1cmp.inconclusive_cells} the shift lies within what the session
effect could produce (Section~\ref{sec:pcheck}), so those cells are not
evidence of a version effect on the GPU comparison. On the baseline suite
the reduction grows at every step in scale on each of these hosts, as the bridge
does\claim{gap-grows-baseline}.
% BEGIN-IF:pcheck
\ifpcheck
These reductions use E1 as measured. On the H100, E1 is inflated by a slower
side-study session (Section~\ref{sec:pcheck}). Corrected for that
(E1/K1, exploratory, point estimates), the H100 reductions are \stat{p2.h100.e1adj_gap.baseline.range} (baseline)
and \stat{p2.h100.e1adj_gap.complex.range} (complex). The adjusted ratio stays below C1 at every scale\claim{h100-adj-shrinks}.
The corrected reductions are:
\stat{p2.e1adj_gap.by_host}.
\fi
% END-IF:pcheck


\begin{table}[t]
\centering
\caption{Pre-registered speedups (C1, C4; paired $t$-intervals, df~$=3$) next
to their \textsc{exploratory} re-baselined counterparts (E1, E4\expl; CPU side
Polars \polarsnew{} from the side study; Welch-type intervals). E1/K1\expl: E1
corrected for the measured session effect (point estimate; Section~\ref{sec:pcheck}).
Cells where the shift from C1 to E1 lies within the session-check range:
\stat{p2.k1cmp.inconclusive_cells}. Above~1: the GPU side is faster. The H100 rows repeat Paper~1's exploratory table for
comparison. \chartnote}
\label{tab:rebase}
\scriptsize
\setlength{\tabcolsep}{2.5pt}
\begin{tabular}{@{}lllccccc@{}}
\toprule
Host & Suite & Scale & C1 & E1\expl & E1/K1\expl & C4 & E4\expl \\
\midrule
% BEGIN-IF:bridge_a100
\ifbridgeaonehundred
A100 & baseline & SF1   & \ratioci{a100.c1.baseline.sf1} & \ratioci{a100.exploratory.e1.baseline.sf1} & \stat{p2.a100.e1adj.baseline.sf1} & \ratioci{a100.c4.baseline.sf1} & \ratioci{a100.exploratory.e4.baseline.sf1} \\
     &          & SF10  & \ratioci{a100.c1.baseline.sf10} & \ratioci{a100.exploratory.e1.baseline.sf10} & \stat{p2.a100.e1adj.baseline.sf10} & \ratioci{a100.c4.baseline.sf10} & \ratioci{a100.exploratory.e4.baseline.sf10} \\
     &          & SF30  & \ratioci{a100.c1.baseline.sf30} & \ratioci{a100.exploratory.e1.baseline.sf30} & \stat{p2.a100.e1adj.baseline.sf30} & \ratioci{a100.c4.baseline.sf30} & \ratioci{a100.exploratory.e4.baseline.sf30} \\
     &          & SF100 & \ratioci{a100.c1.baseline.sf100} & \ratioci{a100.exploratory.e1.baseline.sf100} & \stat{p2.a100.e1adj.baseline.sf100} & \ratioci{a100.c4.baseline.sf100} & \ratioci{a100.exploratory.e4.baseline.sf100} \\
     &          & SF300 & \ratioci{a100.c1.baseline.sf300} & \ratioci{a100.exploratory.e1.baseline.sf300} & \stat{p2.a100.e1adj.baseline.sf300} & \ratioci{a100.c4.baseline.sf300} & \ratioci{a100.exploratory.e4.baseline.sf300} \\
     & complex  & SF30  & \ratioci{a100.c1.complex.sf30} & \ratioci{a100.exploratory.e1.complex.sf30} & \stat{p2.a100.e1adj.complex.sf30} & \ratioci{a100.c4.complex.sf30} & \ratioci{a100.exploratory.e4.complex.sf30} \\
     &          & SF100 & \ratioci{a100.c1.complex.sf100} & \ratioci{a100.exploratory.e1.complex.sf100} & \stat{p2.a100.e1adj.complex.sf100} & \ratioci{a100.c4.complex.sf100} & \ratioci{a100.exploratory.e4.complex.sf100} \\
     &          & SF300 & \ratioci{a100.c1.complex.sf300} & \ratioci{a100.exploratory.e1.complex.sf300} & \stat{p2.a100.e1adj.complex.sf300} & \ratioci{a100.c4.complex.sf300} & \ratioci{a100.exploratory.e4.complex.sf300} \\
\midrule
\fi
% END-IF:bridge_a100
H100 & baseline & SF1   & \ratioci{c1.baseline.sf1} & \ratioci{exploratory.e1.baseline.sf1} & \stat{p2.h100.e1adj.baseline.sf1} & \ratioci{c4.baseline.sf1} & \ratioci{exploratory.e4.baseline.sf1} \\
     &          & SF10  & \ratioci{c1.baseline.sf10} & \ratioci{exploratory.e1.baseline.sf10} & \stat{p2.h100.e1adj.baseline.sf10} & \ratioci{c4.baseline.sf10} & \ratioci{exploratory.e4.baseline.sf10} \\
     &          & SF30  & \ratioci{c1.baseline.sf30} & \ratioci{exploratory.e1.baseline.sf30} & \stat{p2.h100.e1adj.baseline.sf30} & \ratioci{c4.baseline.sf30} & \ratioci{exploratory.e4.baseline.sf30} \\
     &          & SF100 & \ratioci{c1.baseline.sf100} & \ratioci{exploratory.e1.baseline.sf100} & \stat{p2.h100.e1adj.baseline.sf100} & \ratioci{c4.baseline.sf100} & \ratioci{exploratory.e4.baseline.sf100} \\
     &          & SF300 & \ratioci{c1.baseline.sf300} & \ratioci{exploratory.e1.baseline.sf300} & \stat{p2.h100.e1adj.baseline.sf300} & \ratioci{c4.baseline.sf300} & \ratioci{exploratory.e4.baseline.sf300} \\
     & complex  & SF30  & \ratioci{c1.complex.sf30} & \ratioci{exploratory.e1.complex.sf30} & \stat{p2.h100.e1adj.complex.sf30} & \ratioci{c4.complex.sf30} & \ratioci{exploratory.e4.complex.sf30} \\
     &          & SF100 & \ratioci{c1.complex.sf100} & \ratioci{exploratory.e1.complex.sf100} & \stat{p2.h100.e1adj.complex.sf100} & \ratioci{c4.complex.sf100} & \ratioci{exploratory.e4.complex.sf100} \\
     &          & SF300 & \ratioci{c1.complex.sf300} & \ratioci{exploratory.e1.complex.sf300} & \stat{p2.h100.e1adj.complex.sf300} & \ratioci{c4.complex.sf300} & \ratioci{exploratory.e4.complex.sf300} \\
\midrule
GH200 & baseline & SF1   & \ratioci{gh200.c1.baseline.sf1} & \ratioci{gh200.exploratory.e1.baseline.sf1} & \stat{p2.gh200.e1adj.baseline.sf1} & \ratioci{gh200.c4.baseline.sf1} & \ratioci{gh200.exploratory.e4.baseline.sf1} \\
      &          & SF10  & \ratioci{gh200.c1.baseline.sf10} & \ratioci{gh200.exploratory.e1.baseline.sf10} & \stat{p2.gh200.e1adj.baseline.sf10} & \ratioci{gh200.c4.baseline.sf10} & \ratioci{gh200.exploratory.e4.baseline.sf10} \\
      &          & SF30  & \ratioci{gh200.c1.baseline.sf30} & \ratioci{gh200.exploratory.e1.baseline.sf30} & \stat{p2.gh200.e1adj.baseline.sf30} & \ratioci{gh200.c4.baseline.sf30} & \ratioci{gh200.exploratory.e4.baseline.sf30} \\
      &          & SF100 & \ratioci{gh200.c1.baseline.sf100} & \ratioci{gh200.exploratory.e1.baseline.sf100} & \stat{p2.gh200.e1adj.baseline.sf100} & \ratioci{gh200.c4.baseline.sf100} & \ratioci{gh200.exploratory.e4.baseline.sf100} \\
      &          & SF300 & \ratioci{gh200.c1.baseline.sf300} & \ratioci{gh200.exploratory.e1.baseline.sf300} & \stat{p2.gh200.e1adj.baseline.sf300} & \ratioci{gh200.c4.baseline.sf300} & \ratioci{gh200.exploratory.e4.baseline.sf300} \\
      & complex  & SF30  & \ratioci{gh200.c1.complex.sf30} & \ratioci{gh200.exploratory.e1.complex.sf30} & \stat{p2.gh200.e1adj.complex.sf30} & \ratioci{gh200.c4.complex.sf30} & \ratioci{gh200.exploratory.e4.complex.sf30} \\
      &          & SF100 & \ratioci{gh200.c1.complex.sf100} & \ratioci{gh200.exploratory.e1.complex.sf100} & \stat{p2.gh200.e1adj.complex.sf100} & \ratioci{gh200.c4.complex.sf100} & \ratioci{gh200.exploratory.e4.complex.sf100} \\
      &          & SF300 & \ratioci{gh200.c1.complex.sf300} & \ratioci{gh200.exploratory.e1.complex.sf300} & \stat{p2.gh200.e1adj.complex.sf300} & \ratioci{gh200.c4.complex.sf300} & \ratioci{gh200.exploratory.e4.complex.sf300} \\
% BEGIN-IF:bridge_b200
\ifbridgebtwohundred
\midrule
B200 & baseline & SF1   & \ratioci{b200.c1.baseline.sf1} & \ratioci{b200.exploratory.e1.baseline.sf1} & \stat{p2.b200.e1adj.baseline.sf1} & \ratioci{b200.c4.baseline.sf1} & \ratioci{b200.exploratory.e4.baseline.sf1} \\
     &          & SF10  & \ratioci{b200.c1.baseline.sf10} & \ratioci{b200.exploratory.e1.baseline.sf10} & \stat{p2.b200.e1adj.baseline.sf10} & \ratioci{b200.c4.baseline.sf10} & \ratioci{b200.exploratory.e4.baseline.sf10} \\
     &          & SF30  & \ratioci{b200.c1.baseline.sf30} & \ratioci{b200.exploratory.e1.baseline.sf30} & \stat{p2.b200.e1adj.baseline.sf30} & \ratioci{b200.c4.baseline.sf30} & \ratioci{b200.exploratory.e4.baseline.sf30} \\
     &          & SF100 & \ratioci{b200.c1.baseline.sf100} & \ratioci{b200.exploratory.e1.baseline.sf100} & \stat{p2.b200.e1adj.baseline.sf100} & \ratioci{b200.c4.baseline.sf100} & \ratioci{b200.exploratory.e4.baseline.sf100} \\
     &          & SF300 & \ratioci{b200.c1.baseline.sf300} & \ratioci{b200.exploratory.e1.baseline.sf300} & \stat{p2.b200.e1adj.baseline.sf300} & \ratioci{b200.c4.baseline.sf300} & \ratioci{b200.exploratory.e4.baseline.sf300} \\
     & complex  & SF30  & \ratioci{b200.c1.complex.sf30} & \ratioci{b200.exploratory.e1.complex.sf30} & \stat{p2.b200.e1adj.complex.sf30} & \ratioci{b200.c4.complex.sf30} & \ratioci{b200.exploratory.e4.complex.sf30} \\
     &          & SF100 & \ratioci{b200.c1.complex.sf100} & \ratioci{b200.exploratory.e1.complex.sf100} & \stat{p2.b200.e1adj.complex.sf100} & \ratioci{b200.c4.complex.sf100} & \ratioci{b200.exploratory.e4.complex.sf100} \\
     &          & SF300 & \ratioci{b200.c1.complex.sf300} & \ratioci{b200.exploratory.e1.complex.sf300} & \stat{p2.b200.e1adj.complex.sf300} & \ratioci{b200.c4.complex.sf300} & \ratioci{b200.exploratory.e4.complex.sf300} \\
\fi
% END-IF:bridge_b200
\bottomrule
\end{tabular}
\end{table}

\begin{figure}[t]
\centering
\iffigure{f20_baseline_sensitivity}{%
\studyfig{f20_baseline_sensitivity}{\linewidth}}{%
\begin{subfigure}[t]{0.49\linewidth}\centering
\studyfig{f02x_h100}{\linewidth}\caption{H100}\end{subfigure}\hfill
\begin{subfigure}[t]{0.49\linewidth}\centering
\studyfig{f02x_gh200}{\linewidth}\caption{GH200}\end{subfigure}}
\caption{Baseline sensitivity (\textsc{exploratory}, post hoc): the Polars GPU
speedup against Polars \polarsold{} (C1, pre-registered, paired intervals) and
against Polars \polarsnew{} (E1, exploratory, Welch-type intervals), by scale
and suite\iffigure{f20_baseline_sensitivity}{, one column per host with each
host's CPU layout under its name; a host whose side study is not in this
build's inputs shows ``bridge not available''}{; the study's F2x panels also
show C4 and E4}. Above~1: the GPU side is faster.
\chartnote}
\label{fig:sensitivity}
\end{figure}

\paragraph{Where the verdict changes (exploratory).}
A lower speedup matters most when it moves the scale at which the GPU wins.
Table~\ref{tab:crossover} gives, for each host and suite, the first scale at
which each ratio's interval lies entirely above~1\stat{p2.bridge.scope_note}. Changing the CPU baseline from \polarsold{} to
\polarsnew{} moves that scale in these cases: \stat{p2.verdict.moves_list}. It
moves no scale on: \stat{p2.verdict.unchanged_hosts}. In the following cells the
two baselines give opposite verdicts, each with an interval excluding~1: the GPU
engine is faster than Polars \polarsold{} but slower than Polars \polarsnew{}
(\stat{p2.reversal.detail})\claim{reversal-exists}. At those scales on those
hosts, the choice of CPU release, which the GPU package constrains, decides
whether the GPU engine helps.
% BEGIN-IF:pcheck
\ifpcheck
The session correction (E1/K1; point estimates) moves E1's first scale in
these cases: \stat{p2.verdict_adj.moves_list}. The adjusted point estimate at
E1's first scale is: \stat{p2.verdict_adj.at_first_list}.
% BEGIN-IF:e1adj
\ifeadj
With intervals (Table~\ref{tab:e1adj}), the adjusted ratio's first scale
entirely above~1 differs from E1's in these cases:
\stat{p2.verdict_adjci.moves_list}\claim{adj-ci-sustained}. Classifying each
interval as GPU faster, CPU faster or inconclusive, the adjustment changes the
suite-level verdict in \stat{p2.adj.suite_changed_n} of \stat{p2.adj.suite_cells}
cells (E1 and E4: \stat{p2.adj.suite_changed_list}). Per query (E1) it changes
\stat{p2.adj.q_changed_n} of \stat{p2.adj.q_cells}:
\stat{p2.adj.q_changed_list}. Per-query intervals rest on single queries and
four rounds per side, so these changes are fragile, and we draw no conclusion
from them.
\fi
% END-IF:e1adj
% BEGIN-IFNOT:e1adj
\ifeadj\else
Whether the adjusted intervals exclude~1 at those scales is not established in
this build.
\fi
% END-IFNOT:e1adj
\fi
% END-IF:pcheck


\begin{table}[t]
\centering
\caption{First scale at which the 95\% interval lies entirely above~1 (the GPU
side faster), per host and suite. C1/C4: pre-registered contrasts (the C4
column is the pre-registered RQ3 crossover). E1/E4\expl: \textsc{exploratory}. E1/K1\expl: session-adjusted E1, first
scale whose point estimate exceeds~1 and stays above (\textsc{exploratory}).
``---'': no bridge for this host in this build. ``None observed'': no scale
qualifies.}
\label{tab:crossover}
\small
\setlength{\tabcolsep}{4pt}
\begin{tabular}{@{}lcccccccccc@{}}
\toprule
 & \multicolumn{5}{c}{baseline suite} & \multicolumn{5}{c}{complex suite} \\
\cmidrule(lr){2-6}\cmidrule(l){7-11}
Host & C1 & E1\expl & E1/K1\expl & C4 & E4\expl & C1 & E1\expl & E1/K1\expl & C4 & E4\expl \\
\midrule
% BEGIN-IF:bridge_a100
\ifbridgeaonehundred
A100 & \stat{p2.a100.c1.baseline.first_above} & \stat{p2.a100.e1.baseline.first_above} & \stat{p2.a100.e1adj.baseline.first_above_point} & \stat{p2.a100.c4.baseline.first_above} & \stat{p2.a100.e4.baseline.first_above} & \stat{p2.a100.c1.complex.first_above} & \stat{p2.a100.e1.complex.first_above} & \stat{p2.a100.e1adj.complex.first_above_point} & \stat{p2.a100.c4.complex.first_above} & \stat{p2.a100.e4.complex.first_above} \\
\fi
% END-IF:bridge_a100
% BEGIN-IFNOT:bridge_a100
\ifbridgeaonehundred\else
A100 & \stat{p2.a100.c1.baseline.first_above} & --- & --- & \stat{p2.a100.c4.baseline.first_above} & --- & \stat{p2.a100.c1.complex.first_above} & --- & --- & \stat{p2.a100.c4.complex.first_above} & --- \\
\fi
% END-IFNOT:bridge_a100
H100 & \stat{p2.h100.c1.baseline.first_above} & \stat{p2.h100.e1.baseline.first_above} & \stat{p2.h100.e1adj.baseline.first_above_point} & \stat{p2.h100.c4.baseline.first_above} & \stat{p2.h100.e4.baseline.first_above} & \stat{p2.h100.c1.complex.first_above} & \stat{p2.h100.e1.complex.first_above} & \stat{p2.h100.e1adj.complex.first_above_point} & \stat{p2.h100.c4.complex.first_above} & \stat{p2.h100.e4.complex.first_above} \\
GH200 & \stat{p2.gh200.c1.baseline.first_above} & \stat{p2.gh200.e1.baseline.first_above} & \stat{p2.gh200.e1adj.baseline.first_above_point} & \stat{p2.gh200.c4.baseline.first_above} & \stat{p2.gh200.e4.baseline.first_above} & \stat{p2.gh200.c1.complex.first_above} & \stat{p2.gh200.e1.complex.first_above} & \stat{p2.gh200.e1adj.complex.first_above_point} & \stat{p2.gh200.c4.complex.first_above} & \stat{p2.gh200.e4.complex.first_above} \\
% BEGIN-IF:bridge_b200
\ifbridgebtwohundred
B200 & \stat{p2.b200.c1.baseline.first_above} & \stat{p2.b200.e1.baseline.first_above} & \stat{p2.b200.e1adj.baseline.first_above_point} & \stat{p2.b200.c4.baseline.first_above} & \stat{p2.b200.e4.baseline.first_above} & \stat{p2.b200.c1.complex.first_above} & \stat{p2.b200.e1.complex.first_above} & \stat{p2.b200.e1adj.complex.first_above_point} & \stat{p2.b200.c4.complex.first_above} & \stat{p2.b200.e4.complex.first_above} \\
\fi
% END-IF:bridge_b200
% BEGIN-IFNOT:bridge_b200
\ifbridgebtwohundred\else
B200 & \stat{p2.b200.c1.baseline.first_above} & --- & --- & \stat{p2.b200.c4.baseline.first_above} & --- & \stat{p2.b200.c1.complex.first_above} & --- & --- & \stat{p2.b200.c4.complex.first_above} & --- \\
\fi
% END-IFNOT:bridge_b200
\bottomrule
\end{tabular}
\end{table}


% BEGIN-IF:pcheck
\ifpcheck
\subsection{Session bias in E1 (exploratory)}
\label{sec:pcheck}
E1 takes its CPU times from side-study sessions and its GPU times from main
rounds. The change from C1 to E1 therefore mixes two effects: the Polars release,
and any difference between the two kinds of session (time of day, preceding
sessions, file-cache state). Both kinds of session ran Polars \polarsold{}, so
$\mathrm{K1} = t^{\text{side}}_{\text{\polarsold}}/t^{\text{main}}_{\text{\polarsold}}$
measures the session effect on its own.

On the H100, K1 ranged from \stat{p2.h100.pcheck.min} to
\stat{p2.h100.pcheck.max}, and every point estimate was above~1\claim{h100-pcheck-above-1}.
One interval excluded~1
(\stat{p2.h100.pcheck.excluding_1_list}:
\ratiotext{protocol_check.h100.complex.sf100})\claim{h100-pcheck-one-excluded}. On the GH200 no interval excluded~1\claim{gh200-pcheck-none}. For every host:
\stat{p2.pcheck.by_host} (all exploratory, Welch-type).
A side-study session slower by a factor K1 inflates E1 by the same factor.
\textbf{On the H100, E1 may therefore overstate the GPU's advantage over Polars
\polarsnew{} by up to about \stat{p2.h100.pcheck.max_pct}\%.} On the GH200 the bias is negligible. On the A100 the intervals that exclude~1
lie below it\claim{a100-pcheck-below}. There the side-study sessions were
faster, so E1 may \emph{understate} the GPU's advantage over Polars
\polarsnew{}, by up to about \stat{p2.a100.pcheck.min_pct}\%, and the reductions
may overstate the version effect. In every cell the session effect is smaller than the
version effect: K1 lies closer to~1 than the bridge over the GPU-supported
queries\claim{pcheck-small}.

The session-adjusted ratio E1/K1 removes this bias if the session effect is the
same for both Polars releases. The check cannot establish that assumption.
Table~\ref{tab:rebase} gives E1/K1 beside E1, and Section~\ref{sec:res-version} uses it in each statement about E1.

\paragraph{Is the shift from C1 to E1 larger than the session effect?}
For each cell we compare the shift $1-\mathrm{E1}/\mathrm{C1}$ (point estimates)
with the largest distance of the K1 interval from~1. The shift exceeds it in
\stat{p2.k1cmp.summary}. It does not in \stat{p2.k1cmp.inconclusive_all}. In
those cells the re-baselining shift is not distinguishable from the
session-protocol uncertainty, and we do not count them as evidence that the
release changes the GPU comparison. What is inconclusive there is the
comparison of E1 with C1 across the two protocols. The bridge itself is paired
within the side-study rounds, is unaffected, and shows Polars \polarsnew{}
faster in those cells \stat{p2.k1cmp.bridge_phrase_ref}. Every cell where the two
baselines give opposite verdicts has a shift beyond the session-check
range\claim{k1cmp-reversal-exceeds}. This comparison is descriptive and
exploratory.
% BEGIN-IF:e1adj
\ifeadj
Table~\ref{tab:e1adj} gives E1/K1 with intervals, estimated as the bridge over
E1's queries times C1 (both paired within their rounds; Welch--Satterthwaite
df for the combination)\claim{e1adj-matches}.

\begin{table}[t]
\centering
\caption{\textsc{Exploratory}: session-adjusted E1, $\mathrm{E1}/\mathrm{K1}$,
with 95\% intervals, beside the pre-registered C1 and the unadjusted E1. Above~1:
the GPU side is faster. \chartnote}
\label{tab:e1adj}
\scriptsize
\setlength{\tabcolsep}{3pt}
\begin{tabular}{@{}lllccc@{}}
\toprule
Host & Suite & Scale & C1 & E1\expl & E1/K1\expl \\
\midrule
% BEGIN-IF:bridge_a100
\ifbridgeaonehundred
A100 & baseline & SF1 & \ratioci{a100.c1.baseline.sf1} & \ratioci{a100.exploratory.e1.baseline.sf1} & \ratioci{exploratory_adj.e1.a100.baseline.sf1} \\
 &  & SF10 & \ratioci{a100.c1.baseline.sf10} & \ratioci{a100.exploratory.e1.baseline.sf10} & \ratioci{exploratory_adj.e1.a100.baseline.sf10} \\
 &  & SF30 & \ratioci{a100.c1.baseline.sf30} & \ratioci{a100.exploratory.e1.baseline.sf30} & \ratioci{exploratory_adj.e1.a100.baseline.sf30} \\
 &  & SF100 & \ratioci{a100.c1.baseline.sf100} & \ratioci{a100.exploratory.e1.baseline.sf100} & \ratioci{exploratory_adj.e1.a100.baseline.sf100} \\
 &  & SF300 & \ratioci{a100.c1.baseline.sf300} & \ratioci{a100.exploratory.e1.baseline.sf300} & \ratioci{exploratory_adj.e1.a100.baseline.sf300} \\
 & complex & SF30 & \ratioci{a100.c1.complex.sf30} & \ratioci{a100.exploratory.e1.complex.sf30} & \ratioci{exploratory_adj.e1.a100.complex.sf30} \\
 &  & SF100 & \ratioci{a100.c1.complex.sf100} & \ratioci{a100.exploratory.e1.complex.sf100} & \ratioci{exploratory_adj.e1.a100.complex.sf100} \\
 &  & SF300 & \ratioci{a100.c1.complex.sf300} & \ratioci{a100.exploratory.e1.complex.sf300} & \ratioci{exploratory_adj.e1.a100.complex.sf300} \\
\midrule
\fi
% END-IF:bridge_a100
H100 & baseline & SF1   & \ratioci{c1.baseline.sf1} & \ratioci{exploratory.e1.baseline.sf1} & \ratioci{exploratory_adj.e1.h100.baseline.sf1} \\
     &          & SF10  & \ratioci{c1.baseline.sf10} & \ratioci{exploratory.e1.baseline.sf10} & \ratioci{exploratory_adj.e1.h100.baseline.sf10} \\
     &          & SF30  & \ratioci{c1.baseline.sf30} & \ratioci{exploratory.e1.baseline.sf30} & \ratioci{exploratory_adj.e1.h100.baseline.sf30} \\
     &          & SF100 & \ratioci{c1.baseline.sf100} & \ratioci{exploratory.e1.baseline.sf100} & \ratioci{exploratory_adj.e1.h100.baseline.sf100} \\
     &          & SF300 & \ratioci{c1.baseline.sf300} & \ratioci{exploratory.e1.baseline.sf300} & \ratioci{exploratory_adj.e1.h100.baseline.sf300} \\
     & complex  & SF30  & \ratioci{c1.complex.sf30} & \ratioci{exploratory.e1.complex.sf30} & \ratioci{exploratory_adj.e1.h100.complex.sf30} \\
     &          & SF100 & \ratioci{c1.complex.sf100} & \ratioci{exploratory.e1.complex.sf100} & \ratioci{exploratory_adj.e1.h100.complex.sf100} \\
     &          & SF300 & \ratioci{c1.complex.sf300} & \ratioci{exploratory.e1.complex.sf300} & \ratioci{exploratory_adj.e1.h100.complex.sf300} \\
\midrule
GH200 & baseline & SF1   & \ratioci{gh200.c1.baseline.sf1} & \ratioci{gh200.exploratory.e1.baseline.sf1} & \ratioci{exploratory_adj.e1.gh200.baseline.sf1} \\
      &          & SF10  & \ratioci{gh200.c1.baseline.sf10} & \ratioci{gh200.exploratory.e1.baseline.sf10} & \ratioci{exploratory_adj.e1.gh200.baseline.sf10} \\
      &          & SF30  & \ratioci{gh200.c1.baseline.sf30} & \ratioci{gh200.exploratory.e1.baseline.sf30} & \ratioci{exploratory_adj.e1.gh200.baseline.sf30} \\
      &          & SF100 & \ratioci{gh200.c1.baseline.sf100} & \ratioci{gh200.exploratory.e1.baseline.sf100} & \ratioci{exploratory_adj.e1.gh200.baseline.sf100} \\
      &          & SF300 & \ratioci{gh200.c1.baseline.sf300} & \ratioci{gh200.exploratory.e1.baseline.sf300} & \ratioci{exploratory_adj.e1.gh200.baseline.sf300} \\
      & complex  & SF30  & \ratioci{gh200.c1.complex.sf30} & \ratioci{gh200.exploratory.e1.complex.sf30} & \ratioci{exploratory_adj.e1.gh200.complex.sf30} \\
      &          & SF100 & \ratioci{gh200.c1.complex.sf100} & \ratioci{gh200.exploratory.e1.complex.sf100} & \ratioci{exploratory_adj.e1.gh200.complex.sf100} \\
      &          & SF300 & \ratioci{gh200.c1.complex.sf300} & \ratioci{gh200.exploratory.e1.complex.sf300} & \ratioci{exploratory_adj.e1.gh200.complex.sf300} \\
% BEGIN-IF:bridge_b200
\ifbridgebtwohundred
\midrule
B200 & baseline & SF1 & \ratioci{b200.c1.baseline.sf1} & \ratioci{b200.exploratory.e1.baseline.sf1} & \ratioci{exploratory_adj.e1.b200.baseline.sf1} \\
 &  & SF10 & \ratioci{b200.c1.baseline.sf10} & \ratioci{b200.exploratory.e1.baseline.sf10} & \ratioci{exploratory_adj.e1.b200.baseline.sf10} \\
 &  & SF30 & \ratioci{b200.c1.baseline.sf30} & \ratioci{b200.exploratory.e1.baseline.sf30} & \ratioci{exploratory_adj.e1.b200.baseline.sf30} \\
 &  & SF100 & \ratioci{b200.c1.baseline.sf100} & \ratioci{b200.exploratory.e1.baseline.sf100} & \ratioci{exploratory_adj.e1.b200.baseline.sf100} \\
 &  & SF300 & \ratioci{b200.c1.baseline.sf300} & \ratioci{b200.exploratory.e1.baseline.sf300} & \ratioci{exploratory_adj.e1.b200.baseline.sf300} \\
 & complex & SF30 & \ratioci{b200.c1.complex.sf30} & \ratioci{b200.exploratory.e1.complex.sf30} & \ratioci{exploratory_adj.e1.b200.complex.sf30} \\
 &  & SF100 & \ratioci{b200.c1.complex.sf100} & \ratioci{b200.exploratory.e1.complex.sf100} & \ratioci{exploratory_adj.e1.b200.complex.sf100} \\
 &  & SF300 & \ratioci{b200.c1.complex.sf300} & \ratioci{b200.exploratory.e1.complex.sf300} & \ratioci{exploratory_adj.e1.b200.complex.sf300} \\
\fi
% END-IF:bridge_b200
\bottomrule
\end{tabular}
\end{table}
\fi
% END-IF:e1adj
\fi
% END-IF:pcheck

\subsection{Equal threads, unequal cores}
\label{sec:res-cores}

\paragraph{Topology (recorded).} Table~\ref{tab:topology} lists what each guest
reported. On every host the engines ran with \stat{host.cpus_arg} threads\claim{same-threads}. On the H100 and B200 those threads match the vCPU count and
occupy \stat{p2.h100.cores_used} physical cores, two hardware threads per core\claim{h100-smt}. On the A100 and GH200 they occupy
\stat{p2.a100.cores_used} cores, one thread per core, and the remaining vCPUs
are not used by the worker pools\claim{cores-differ}\claim{cores-pairs}. The design meant ``same threads on every host''
as a control (amendment A3). For CPU engines it does not hold the CPU resources
constant: the H100 and B200 baselines run on half as many cores. Because
threads are not pinned, nothing stops helper threads (I/O, garbage collection,
the operating system) from using the idle cores on the A100 and GH200. On the
H100 and B200 the worker threads already equal the vCPU count.

\begin{table}[t]
\centering
\caption{CPU as reported by each guest (\texttt{lscpu}, recorded by the runner's standard inventory; analysed post hoc). Physical cores are as the guest reports them; a hypervisor can expose
a topology that differs from the physical machine.}
\label{tab:topology}
\footnotesize
\begin{tabular}{@{}>{\raggedright\arraybackslash}p{3.2cm}*{4}{>{\raggedright\arraybackslash}p{2.75cm}}@{}}
\toprule
 & A100 & H100 & GH200 & B200 \\
\midrule
CPU model & \stat{a100.host.cpu_model} & \stat{host.cpu_model} & \stat{gh200.host.cpu_model} & \stat{b200.host.cpu_model} \\
Architecture & \stat{a100.host.arch} & \stat{host.arch} & \stat{gh200.host.arch} & \stat{b200.host.arch} \\
Hypervisor line & \stat{p2.a100.hypervisor_text} & \stat{p2.h100.hypervisor_text} & \stat{p2.gh200.hypervisor_text} & \stat{p2.b200.hypervisor_text} \\
vCPUs & \stat{a100.host.vcpus} & \stat{host.vcpus} & \stat{gh200.host.vcpus} & \stat{b200.host.vcpus} \\
Hardware threads per core & \stat{a100.host.smt_threads_per_core} & \stat{host.smt_threads_per_core} & \stat{gh200.host.smt_threads_per_core} & \stat{b200.host.smt_threads_per_core} \\
Physical cores & \stat{a100.host.physical_cores} & \stat{host.physical_cores} & \stat{gh200.host.physical_cores} & \stat{b200.host.physical_cores} \\
NUMA nodes (with CPUs) & \stat{a100.host.numa_nodes} (\stat{a100.host.numa_nodes_with_cpus}) & \stat{host.numa_nodes} (\stat{host.numa_nodes_with_cpus}) & \stat{gh200.host.numa_nodes} (\stat{gh200.host.numa_nodes_with_cpus}) & \stat{b200.host.numa_nodes} (\stat{b200.host.numa_nodes_with_cpus}) \\
Engine threads (\texttt{-{}-cpus}) & \stat{a100.host.cpus_arg} & \stat{host.cpus_arg} & \stat{gh200.host.cpus_arg} & \stat{b200.host.cpus_arg} \\
Physical cores they occupy & \stat{p2.a100.cores_used} & \stat{p2.h100.cores_used} & \stat{p2.gh200.cores_used} & \stat{p2.b200.cores_used} \\
Share of vCPUs used & \stat{a100.host.cpus_arg_share_of_vcpus} & \stat{host.cpus_arg_share_of_vcpus} & \stat{gh200.host.cpus_arg_share_of_vcpus} & \stat{b200.host.cpus_arg_share_of_vcpus} \\
\bottomrule
\end{tabular}
\end{table}

\paragraph{CPU-engine times across hosts (descriptive).}
Figure~\ref{fig:ladder} shows each engine's time on each host relative to the
H100, and Table~\ref{tab:xhost} gives the values at SF300. Polars CPU on the
GH200, whose threads sit on \stat{p2.gh200.cores_used} cores, took between
\stat{xhost.gh200.polars.overall.min} and \stat{xhost.gh200.polars.overall.max}
of the H100's time, with every interval below~1\claim{gh200-cpu-faster}.
Spark CPU on the GH200 took between \stat{xhost.gh200.spark_cpu.overall.min}
and \stat{xhost.gh200.spark_cpu.overall.max}. Core count alone does not explain
these numbers. Polars CPU on the A100 also runs on
\stat{p2.a100.cores_used} cores, of a different CPU model, and took between
\stat{xhost.a100.polars.overall.min} and \stat{xhost.a100.polars.overall.max}
of the H100's time, with every interval above~1\claim{a100-cpu-slower}. The
B200 matches the H100's topology. There Polars CPU took between
\stat{xhost.b200.polars.overall.min} and \stat{xhost.b200.polars.overall.max}
of the H100's time. The baseline therefore depends on the whole host CPU,
including core count, microarchitecture and memory system, and the shared
thread count equalizes none of them.

\begin{figure}[t]
\centering
\studyfig{f11_gpu_ladder}{\linewidth}
\caption{Suite-level time of each engine on each host relative to the H100
(host/H100; geometric mean of round medians; 95\% Welch-type intervals;
pre-registered as a descriptive secondary analysis, amendments A2/A3). The
CPU-engine rows show the CPU side of each speedup. Figure from the study's
cross-host analysis (F11). \chartnote}
\label{fig:ladder}
\end{figure}

\begin{table}[t]
\centering
\caption{SF300 cross-host time ratios $X = t_{\text{host}}/t_{\text{H100}}$ per
engine (descriptive; Welch-type 95\% intervals) and the host's speedups C1 and
C2 (pre-registered; paired $t$-intervals). By Eq.~\eqref{eq:factor},
$C^{h}/C^{\text{H100}} = X_{\text{cpu}}/X_{\text{gpu}}$ for the point estimates
except C1 on the complex suite (marked $^\dagger$, where $X_{\text{cpu}}$
includes two queries outside C1). The last two columns
(\textsc{exploratory}\expl) give $C^{h}/C^{\text{H100}}$. \chartnote}
\label{tab:xhost}
\scriptsize
\setlength{\tabcolsep}{2pt}
\resizebox{\linewidth}{!}{%
\begin{tabular}{@{}llcccccccc@{}}
\toprule
Suite & Host & cores & $X$ Polars CPU & $X$ Polars GPU & C1 & $X$ Spark CPU & $X$ Spark GPU & C2 & C1, C2 ratio\expl \\
\midrule
baseline & A100  & \stat{p2.a100.cores_used} & \ratioci{xhost.a100.polars.baseline.sf300} & \ratioci{xhost.a100.polars_gpu.baseline.sf300} & \ratioci{a100.c1.baseline.sf300} & \ratioci{xhost.a100.spark_cpu.baseline.sf300} & \ratioci{xhost.a100.spark_gpu.baseline.sf300} & \ratioci{a100.c2.baseline.sf300} & \stat{p2.a100.c1.baseline.sf300.speedup_ratio}, \stat{p2.a100.c2.baseline.sf300.speedup_ratio} \\
         & H100  & \stat{p2.h100.cores_used} & ref. & ref. & \ratioci{c1.baseline.sf300} & ref. & ref. & \ratioci{c2.baseline.sf300} & ref. \\
         & GH200 & \stat{p2.gh200.cores_used} & \ratioci{xhost.gh200.polars.baseline.sf300} & \ratioci{xhost.gh200.polars_gpu.baseline.sf300} & \ratioci{gh200.c1.baseline.sf300} & \ratioci{xhost.gh200.spark_cpu.baseline.sf300} & \ratioci{xhost.gh200.spark_gpu.baseline.sf300} & \ratioci{gh200.c2.baseline.sf300} & \stat{p2.gh200.c1.baseline.sf300.speedup_ratio}, \stat{p2.gh200.c2.baseline.sf300.speedup_ratio} \\
         & B200  & \stat{p2.b200.cores_used} & \ratioci{xhost.b200.polars.baseline.sf300} & \ratioci{xhost.b200.polars_gpu.baseline.sf300} & \ratioci{b200.c1.baseline.sf300} & \ratioci{xhost.b200.spark_cpu.baseline.sf300} & \ratioci{xhost.b200.spark_gpu.baseline.sf300} & \ratioci{b200.c2.baseline.sf300} & \stat{p2.b200.c1.baseline.sf300.speedup_ratio}, \stat{p2.b200.c2.baseline.sf300.speedup_ratio} \\
\midrule
complex  & A100  & \stat{p2.a100.cores_used} & \ratioci{xhost.a100.polars.complex.sf300}$^\dagger$ & \ratioci{xhost.a100.polars_gpu.complex.sf300} & \ratioci{a100.c1.complex.sf300} & \ratioci{xhost.a100.spark_cpu.complex.sf300} & \ratioci{xhost.a100.spark_gpu.complex.sf300} & \ratioci{a100.c2.complex.sf300} & \stat{p2.a100.c1.complex.sf300.speedup_ratio}, \stat{p2.a100.c2.complex.sf300.speedup_ratio} \\
         & H100  & \stat{p2.h100.cores_used} & ref. & ref. & \ratioci{c1.complex.sf300} & ref. & ref. & \ratioci{c2.complex.sf300} & ref. \\
         & GH200 & \stat{p2.gh200.cores_used} & \ratioci{xhost.gh200.polars.complex.sf300}$^\dagger$ & \ratioci{xhost.gh200.polars_gpu.complex.sf300} & \ratioci{gh200.c1.complex.sf300} & \ratioci{xhost.gh200.spark_cpu.complex.sf300} & \ratioci{xhost.gh200.spark_gpu.complex.sf300} & \ratioci{gh200.c2.complex.sf300} & \stat{p2.gh200.c1.complex.sf300.speedup_ratio}, \stat{p2.gh200.c2.complex.sf300.speedup_ratio} \\
         & B200  & \stat{p2.b200.cores_used} & \ratioci{xhost.b200.polars.complex.sf300}$^\dagger$ & \ratioci{xhost.b200.polars_gpu.complex.sf300} & \ratioci{b200.c1.complex.sf300} & \ratioci{xhost.b200.spark_cpu.complex.sf300} & \ratioci{xhost.b200.spark_gpu.complex.sf300} & \ratioci{b200.c2.complex.sf300} & \stat{p2.b200.c1.complex.sf300.speedup_ratio}, \stat{p2.b200.c2.complex.sf300.speedup_ratio} \\
\bottomrule
\end{tabular}}
\end{table}

\paragraph{Factorizing the speedup differences (exploratory).}
Equation~\eqref{eq:factor} splits each host's speedup difference into the two
rows of Table~\ref{tab:xhost}. On the GH200 the Polars GPU speedup was lower
than on the H100 at every scale\claim{gh200-c1-lower}. At SF300 on the
baseline suite it was \ratiotext{gh200.c1.baseline.sf300} against
\ratiotext{c1.baseline.sf300}, a ratio of
\stat{p2.gh200.c1.baseline.sf300.speedup_ratio}. Most of that comes from the CPU
side. In every cell, the GH200's Polars CPU ratio lies further from~1
than its Polars GPU ratio\claim{gh200-cpu-side}. At SF300 they are
\ratiotext{xhost.gh200.polars.baseline.sf300} for the CPU and
\ratiotext{xhost.gh200.polars_gpu.baseline.sf300} for the GPU. Over all cells, the GH200's Polars GPU engine took between
\stat{xhost.gh200.polars_gpu.overall.min} and
\stat{xhost.gh200.polars_gpu.overall.max} of the H100's time. Its speedup is
smaller mainly because its CPU baseline is faster, not because its GPU engine
is slower. The A100 shows the
opposite pattern. Its Polars GPU speedup is also lower than the H100's
(\stat{p2.a100.c1.baseline.sf300.speedup_ratio} of it at SF300 on the baseline
suite), but there the GPU side dominates in every cell\claim{a100-gpu-side}. The A100's GPU engine was slower
(\ratiotext{xhost.a100.polars_gpu.baseline.sf300} of the H100's time), and so
was its CPU engine (\ratiotext{xhost.a100.polars.baseline.sf300}). The
B200 shares the H100's reported topology. Its speedup at SF300 on the baseline
suite was \stat{p2.b200.c1.baseline.sf300.speedup_ratio} of the H100's, with a
CPU-side ratio of \ratiotext{xhost.b200.polars.baseline.sf300} and a GPU-side
ratio of \ratiotext{xhost.b200.polars_gpu.baseline.sf300}.
% BEGIN-IF:decomp
\ifdecomp
Figure~\ref{fig:decomp} shows the factorization for every host, contrast,
suite and scale. It includes C1 on the complex suite, with the CPU side
restricted to the GPU-supported queries, and intervals for the speedup ratio
itself. At SF300 on the baseline suite the GH200's ratio was
\ratiotext{decomp.gh200.c1.baseline.sf300.speedup_ratio}.

\begin{figure}[t]
\centering
\studyfig{f21_speedup_decomposition}{\linewidth}
\caption{\textsc{Exploratory} factorization of cross-host speedup differences
(Eq.~\eqref{eq:factor}): for each host relative to the H100, the speedup ratio
$C^{h}/C^{\text{H100}}$ (filled black cross; Welch-type interval across hosts)
and the CPU and GPU engines' time ratios (open markers in each engine's own
marker and colour, as in the other figures), over the same queries, for C1
(Polars) and C2 (Spark). Each host's CPU layout is printed under its name. A CPU-side value below~1 (the
host's CPU engine faster) lowers the host's speedup. Descriptive; CPU model,
architecture and topology change together across hosts. \chartnote}
\label{fig:decomp}
\end{figure}
\fi
% END-IF:decomp

These comparisons are descriptive. Moving from the H100 to the GH200 changes
the core count, the SMT setting, the CPU vendor and microarchitecture, the
instruction set (x86\_64 to aarch64), the memory system and the platform
builds of every package all at once. The data cannot say which of these makes
the GH200's CPU engines faster. They do establish that a
\emph{same-thread-count} control leaves the CPU baseline free to vary across
hosts by a factor that is comparable to the GPU-side differences, and
sometimes larger.

\subsection{Case: the earlier A100 study}
\label{sec:res-prior}

Before the design was locked, the same project ran an A100 study
(\stat{a100_h100.study.label}). It was not pre-registered: it had no locked
design and a single session per cell. We use it only as an illustrative case. Its saved evidence lets us see how baseline
choices carry into a later comparison. The analysis lists its differences from
the present campaigns as follows: \stat{a100_h100.study.difference_headline}
In particular, it ran Polars \stat{a100_h100.study.polars_version} with
\stat{a100_h100.study.cpu_threads} threads, against Polars
\stat{a100_h100.study.polars_version_h100} with
\stat{a100_h100.study.cpu_threads_h100} threads here\claim{prior-threads}, and
had \stat{a100_h100.study.sessions_per_cell} session per cell against
\stat{a100_h100.study.rounds_h100} rounds. It had no Polars GPU arm.

So no Polars CPU or speedup number carries over from that study to this one.
Its CPU engine ran a different release with more threads. Its A100 instance
had \stat{a100_h100.study.vcpus} vCPUs, and its CPU model and topology were
not recorded.
The only engine the analysis compares across the two studies is Spark GPU
(Table~\ref{tab:prior}). Two estimates of the same quantity, Spark GPU time on
an A100 relative to the H100, come from the earlier study and from this
study's A100 campaign. The two A100 instances may differ too, so their difference
has no single cause. It shows how far apart two careful measurements of the
same cross-host number can lie when protocol, threads and instance all change. Note that the earlier
study's interval reflects only the H100's round-to-round variation, because
that study has a single session per cell.

\begin{table}[t]
\centering
\caption{Spark GPU time on an A100 relative to the H100, from two sources:
the earlier A100 study (run before the design was locked; not pre-registered: no locked design and a single session per cell; an illustrative case only) (\stat{a100_h100.study.cpu_threads}
threads, one session per cell; the interval reflects only the H100's
variation), and this study's A100 campaign (\stat{a100.host.cpus_arg} threads,
four rounds; Welch-type interval). Above~1: the H100 host was faster.
Descriptive. \chartnote}
\label{tab:prior}
\small
\begin{tabular}{@{}llcc@{}}
\toprule
Suite & Scale & earlier A100 study / H100 & A100 campaign / H100 \\
\midrule
baseline & SF1   & \ratioci{a100_h100.spark_gpu.baseline.sf1} & \ratioci{xhost.a100.spark_gpu.baseline.sf1} \\
         & SF10  & \ratioci{a100_h100.spark_gpu.baseline.sf10} & \ratioci{xhost.a100.spark_gpu.baseline.sf10} \\
         & SF30  & \ratioci{a100_h100.spark_gpu.baseline.sf30} & \ratioci{xhost.a100.spark_gpu.baseline.sf30} \\
         & SF100 & \ratioci{a100_h100.spark_gpu.baseline.sf100} & \ratioci{xhost.a100.spark_gpu.baseline.sf100} \\
         & SF300 & \ratioci{a100_h100.spark_gpu.baseline.sf300} & \ratioci{xhost.a100.spark_gpu.baseline.sf300} \\
complex  & SF30  & \ratioci{a100_h100.spark_gpu.complex.sf30} & \ratioci{xhost.a100.spark_gpu.complex.sf30} \\
         & SF100 & \ratioci{a100_h100.spark_gpu.complex.sf100} & \ratioci{xhost.a100.spark_gpu.complex.sf100} \\
         & SF300 & \ratioci{a100_h100.spark_gpu.complex.sf300} & \ratioci{xhost.a100.spark_gpu.complex.sf300} \\
\bottomrule
\end{tabular}
\end{table}

% TODO(results): one sentence on Table tab:prior once read with the author
% (e.g. whether the two sources agree within their intervals at most scales);
% keep it descriptive -- the A100 instances may also differ.

% ======================================================================
\section{Discussion}
\label{sec:discussion}

\paragraph{Two questions, two baselines.} The pre-registered C1 answers one
question: how much faster does this Polars release run a query on the GPU
engine than on its own CPU engine? For that question, sharing one Polars install
is correct, since the two sides differ only in the engine. A user deciding
whether to buy or rent a GPU asks a different question: how much faster is the
GPU than the best CPU setup available? For that question the baseline is the
newest CPU release on the CPU the user would otherwise have. E1 approximates
the newer-release part of the second question on our hosts. On the H100 at
SF300 the two answers differ by \stat{p2.h100.gap.baseline.sf300}\% (baseline
suite) and \stat{p2.h100.gap.complex.sf300}\% (complex suite). After the
session correction they differ by \stat{p2.h100.e1adj_gap.baseline.sf300}\% and
\stat{p2.h100.e1adj_gap.complex.sf300}\%. At intermediate
scales they can even point in opposite directions
(Table~\ref{tab:crossover}). Neither number is wrong.
They are wrong only when reported without saying which question they answer.

\paragraph{Lock-in is structural.} The constraint comes from how the software
is released. The GPU engine is tied to a range of the CPU library's versions
and follows them with a delay~\citep{pypi_cudfpolars_26102,pypi_polars}. During
such a delay, any benchmark of the GPU engine has to use an older CPU release.
A benchmark that does not report a bridge to the current release leaves its
readers unable to tell how much of the speedup is a version gap. The size of
the gap here is specific to Polars \polarsold{} and \polarsnew{} on this
workload; at another moment in the release cycle it may be zero or larger.

\paragraph{Thread counts are not resources.} Fixing the thread count across
hosts is a natural control and was pre-registered as one (A3). It equalizes the
size of the engines' worker pools. It does not equalize cores, SMT, caches,
memory bandwidth, or idle capacity for helper threads, and in cloud guests the
topology is whatever the hypervisor exposes. Cross-host comparisons of GPU
speedups therefore need the CPU-side ratio next to them, as in
Table~\ref{tab:xhost}. Equation~\eqref{eq:factor} makes the split mechanical.
Without it, the GH200's lower Polars GPU speedup could be misread as a weaker
GPU. The cross-host ratios show that most of the difference lies on the CPU
side.

\paragraph{The query set is part of the baseline.} The complex-suite bridge
over all six queries (\stat{p2.bridge.complex.min}--\stat{p2.bridge.complex.max})
is pulled down by one query the GPU engine cannot run. Over the queries that
enter the GPU contrasts it is \stat{p2.bridge_gpuq.complex.min}--\stat{p2.bridge_gpuq.complex.max}.
A version or hardware effect quoted for a CPU baseline should use the same
queries as the speedup it is meant to adjust.

\subsection{A reporting checklist}
\label{sec:checklist}

Table~\ref{tab:checklist} turns these observations into items for reporting
CPU-versus-GPU dataframe benchmarks. Each row also records how the present
study meets it. Items 1--3 concern the version, 4--6 the CPU, and 7--10 the
analysis.

\begin{table}[t]
\centering
\caption{Reporting checklist for CPU-versus-GPU dataframe benchmarks, and how
the source study meets each item.}
\label{tab:checklist}
\footnotesize
\begin{tabular}{@{}r>{\raggedright\arraybackslash}p{6.3cm}>{\raggedright\arraybackslash}p{7.2cm}@{}}
\toprule
\# & Report & This study \\
\midrule
1 & The CPU library version in the baseline, and the newest release at the time of measurement & \polarsold{} in the baseline; \polarsnew{} measured in the bridge; a newer release (2.1.0) existed and was not measured~\citep{pypi_polars} \\
2 & Why the baseline uses that version (e.g.\ the GPU package's supported range) & cudf-polars 26.10.2 requires Polars ${<}1.45$~\citep{pypi_cudfpolars_26102} \\
3 & A bridge from the baseline release to the newest, over the same queries as the speedup & bridge pre-registered (RQ5); restriction to GPU queries post hoc (point estimates) \\
4 & CPU model, architecture, vCPUs, hardware threads per core, physical cores, NUMA nodes with CPUs, hypervisor; say whether guest-reported & recorded by the standard inventory (\texttt{lscpu}); guest-reported; analysed post hoc (Table~\ref{tab:topology}) \\
5 & Engine thread counts \emph{and} the physical cores they occupy; whether threads are pinned; what other threads may run & \stat{host.cpus_arg} threads on \stat{p2.h100.cores_used} to \stat{p2.a100.cores_used} cores depending on host; not pinned; helper threads not capped \\
6 & For cross-host comparisons: the CPU-engine time ratio beside every speedup comparison (Eq.~\ref{eq:factor}) & cross-host ratios pre-registered as descriptive; factorization post hoc \\
7 & Which contrasts, baselines and crossover rules were fixed before data collection & locked design with dated amendments, commit-timestamped \\
8 & Sensitivity analyses with alternative baselines, labelled exploratory, with how their intervals were formed & E1/E4 with Welch-type intervals, labelled exploratory \\
9 & Uncertainty from independent replications (rounds), not only repeated trials within a session & \stat{meta.rounds_completed} rounds per host, Latin-square order \\
10 & Queries excluded from a speedup and why (coverage), so the baseline can be restricted to the same set & pre-declared and observed coverage gaps; excluded from contrasts \\
\bottomrule
\end{tabular}
\end{table}

% ======================================================================
\section{Threats to Validity}
\label{sec:threats}

\paragraph{Post-hoc analysis.} E1, E4, the factorization, the first-above-1
scales for E1/E4, the GPU-query bridge and the session check were chosen after
the bridge and cross-host results were known. Table~\ref{tab:status} lists every post-hoc analysis this paper
uses; the choice of which to run was data-dependent. No
multiplicity correction is applied, and none of these results is confirmatory.

\paragraph{Unpaired sessions.} E1 and E4 combine side-study CPU sessions with
main-round GPU sessions.
% BEGIN-IF:pcheck
\ifpcheck
Section~\ref{sec:pcheck} checks the session effect for Polars \polarsold{}. It
cannot rule out a session effect that differs between Polars releases.
\fi
% END-IF:pcheck
% BEGIN-IFNOT:pcheck
\ifpcheck\else
We did not separately estimate the difference between side-study and
main-round sessions. It enters E1 and E4 alongside the release effect.
\fi
% END-IFNOT:pcheck
The Welch-type intervals treat the rounds on the two sides as independent
samples. With four rounds per side and small within-host variability, they can
be narrow. They do not cover any systematic difference between the session
types.

\paragraph{Guest-reported topology.} Every topology value comes from the
guest. On a shared cloud host, the physical cores behind the vCPUs, whether
sibling hardware threads are really siblings, and contention from other
tenants are all invisible to us.

\paragraph{No within-host topology variation.} No host varied the thread count,
SMT use or pinning. The topology therefore co-varies with everything else that
differs between hosts, and we can only show that the CPU side differs as the
topology would suggest. We cannot show that the topology causes it. A
within-host sweep (threads, SMT on/off, pinning) would separate the factors.

\paragraph{Release-specific results.} The bridge measures one pair of releases
on one workload. Polars~2.1.0 was not measured. The bridge does not generalize
to other releases, to other dataframe engines, or to a workload with different
operators. The \q{cx\_window} result shows how much a single operator can
contribute.

\paragraph{Scope.} The workload is TPC-H-derived and not comparable to TPC-H
results (Section~\ref{sec:scope}). The hosts are single cloud instances, one
per GPU type.

% ======================================================================
\section{Reproducibility}
\label{sec:repro}

This paper uses only the study's analysis outputs: each host's
\texttt{stats.json} and the cross-host files.
% BEGIN-IF:pcheck
\ifpcheck
It also uses a Paper 2 analysis file holding the session check and the
factorization intervals.
\fi
% END-IF:pcheck
The build records the SHA-256 of every input in \texttt{generated/stats.tex}
and refuses to build if any number the text uses is missing. It also fails if
any statement the text makes in words (``at every scale'', ``every interval
below~1'') no longer holds for the inputs. Raw data, the analysis code and the
design document with its amendments are in the study repository
\url{https://github.com/ejiod2025/gpu-spark-polars-study}.
% TODO(author): archive DOI (e.g. Zenodo) for the repository snapshot.

% ======================================================================
\section{Conclusion}
\label{sec:conclusion}

The same GPU engine, on the same host and data, showed GPU speedups that
differed by up to \stat{p2.gap.max}\% depending on which Polars release formed
the baseline (exploratory; \stat{p2.e1adj_gap.max}\% after correcting for the
measured session bias). \stat{p2.reversal.phrase_cap} the two baselines disagreed on whether the GPU
helped at all\claim{reversal-exists}. Across hosts, holding the thread count fixed left the CPU
baseline free to vary. On the GH200 that variation, rather than the GPU,
accounts for most of the lower speedup (exploratory, descriptive). Neither effect is visible in a headline speedup.
Both are visible once the baseline's release, the newest release and the CPU
topology are reported, and the CPU-side ratio is given next to any cross-host
comparison. Our recommendation is to report those items and to label every
alternative baseline as what it is: in this study, an exploratory sensitivity
analysis.

\section*{Acknowledgements}
\todoplaceholder{acknowledgements; funding; compute provider credit if any}

\bibliographystyle{plainnat}
\bibliography{references}

\end{document}
